% ==============================================================================
% Apêndice A: Declaração de uso de inteligência artificial
% Exigida pelo Art. 19 da Portaria Normativa nº 310/2026 da Ufes.
% Se você não usou ferramentas de IA, apague este arquivo, a linha que o
% inclui no arquivo principal e a Seção 6.5 da Implementação.
% ==============================================================================
\chapter{Declaração de uso de inteligência artificial}
\label{ap:ia}

\begin{instrucao}
\orientacao{Por que um apêndice?}
O Art. 19 da Portaria Normativa nº 310/2026 da Ufes \cite{ufes2026portaria}
admite que a declaração fique em seção ou apêndice específico. O apêndice foi
escolhido porque a NBR 14724 não prevê um elemento pré-textual para esse fim e
porque o detalhamento das interações exigido pelo § 2º pode ser longo. Por ser
o primeiro apêndice, ele aparece logo após as Referências e é fácil de
localizar. Na Implementação, uma seção curta cita a ferramenta e remete a este
apêndice.

\orientacao{Como preencher}
\begin{itemize}
  \item No \cref{qua:ia:declaracao}, liste cada ferramenta usada, com a versão
        ou o modelo, a finalidade e a forma de contribuição, indicando as
        partes do trabalho afetadas.
  \item Na seção seguinte, detalhe as interações, informando o que foi pedido,
        por que, como o resultado foi verificado e o que foi aproveitado ou
        descartado.
        Não é preciso transcrever todas as conversas. Descreva os tipos de
        interação com exemplos representativos.
  \item Cite cada ferramenta e inclua-a no arquivo \verb|referencias.bib|
        (veja as entradas \verb|openai2026chatgpt| e \verb|github2026copilot|).
  \item A declaração é escrita na primeira pessoa do singular (``Declaro
        que\ldots''), por ser um compromisso pessoal do autor. O detalhamento
        das interações volta à terceira pessoa, como o restante do trabalho.
  \item Seja específico e honesto. Usos ``meramente acessórios'', que a
        portaria dispensa de declaração, serão definidos pela CPGIA. Na dúvida,
        declare.
\end{itemize}
Substitua o texto de exemplo abaixo pelas ferramentas, versões e usos reais
do seu trabalho. Quando a CPGIA publicar o modelo oficial de declaração,
adapte este apêndice a ele.
\end{instrucao}

% ------------------------------------------------------------------------------
\section{Declaração}

Declaro, em cumprimento ao Art. 19 da Portaria Normativa nº 310/2026 da
\imprimiruniversidade{} \cite{ufes2026portaria}, que utilizei ferramentas de
\acl{ia} na elaboração deste trabalho, nos termos descritos no
\cref{qua:ia:declaracao}. Declaro também que revisei todo o conteúdo gerado
com o auxílio dessas ferramentas, que assumo integral responsabilidade pela
precisão e pela originalidade do texto, do código e dos resultados
apresentados e que nenhuma ferramenta foi tratada como autora ou coautora do
trabalho.

\begin{quadro}[htbp]
  \centering
  \caption{Ferramentas de \acs{ia} utilizadas no trabalho}
  \label{qua:ia:declaracao}
  \small
  \begin{tabularx}{\textwidth}{|>{\raggedright\arraybackslash}p{2.8cm}|>{\raggedright\arraybackslash}X|>{\raggedright\arraybackslash}X|>{\raggedright\arraybackslash}p{2.6cm}|}
    \hline
    \textbf{Ferramenta e versão} & \textbf{Finalidade} &
      \textbf{Forma de contribuição} & \textbf{Partes afetadas} \\
    \hline
    GitHub Copilot, extensão do Visual Studio Code, modelo padrão do plano
      para estudantes \cite{github2026copilot} &
      Apoio à programação &
      Sugestões de complemento de código em funções de formatação de
      relatórios e em testes, revisadas e ajustadas pelo autor &
      Módulo \texttt{relatorios} e testes automatizados
      (\cref{sec:implementacao:organizacao,sec:testes:automatizados}) \\
    \hline
  \end{tabularx}
  \fonteautor
\end{quadro}

% ------------------------------------------------------------------------------
\section{Detalhamento das interações}

\textbf{Complemento de código.} A extensão ficou ativa no editor durante a
implementação do módulo \texttt{relatorios} e dos testes automatizados. Ela
sugeria complementos de linhas e de funções a partir do código já escrito e
dos comentários. A motivação foi reduzir o tempo gasto com código repetitivo,
como a formatação de datas e a montagem de tabelas dos relatórios. As
sugestões foram aceitas apenas quando o autor entendia o código proposto.
Todas foram revisadas, várias foram modificadas e nenhuma foi incorporada sem
teste automatizado correspondente.

\textbf{O que não foi feito com a ferramenta.} A especificação de requisitos,
o projeto da arquitetura, o modelo de dados, a regra de tratamento de
conflitos de reserva (\cref{cod:conflito}) e o texto deste trabalho foram
elaborados sem o uso de ferramentas de \acs{ia}.

\textbf{Configuração.} A extensão foi usada com as configurações padrão,
entre agosto e outubro de 2026, com a opção de compartilhamento de trechos de
código para treinamento desativada. Nenhum dado pessoal, credencial ou dado
dos participantes da avaliação foi fornecido à ferramenta.
